\section{An admissible target with no coefficient-linear tail of primitives}
\label{sec:136-native-tail-nosplit}

\paragraph{Scope and archive locator.}
The vertical and horizontal conditions are those of cote 136,
physical pp.~117 and 119, footers 116 and 118, scans 117 and 119.
The common-source question is on physical p.~121, footer 120,
scan 121. This module tests the proposed replacement of a free
coefficient-tail resolution by a single nonflat tail module of
primitives. It constructs an actual integral admissible target,
and actual guarded source data, for which that particular lift is
impossible. It does not obstruct arbitrary admissible common sources.

\paragraph{ARCH-136-TAIL-HOM-VANISHING.}
Let \(R=\mathbb Z_{(p)}\), \(S=\Gamma_R(T)\), and
\(I_M=\bigoplus_{w\geq M}R T^{[w]}\), where \(M\geq1\).
Write \(S e_s\) for a free graded module on a generator of weight
\(s\). Then
\[
 \operatorname{Hom}^{0}_{S,\mathrm{gr}}(I_M,S e_s)=0
                 \quad(s>0),\qquad
 \operatorname{Hom}^{0}_{S,\mathrm{gr}}(I_M,S)=R.
                                                        \tag{136.NT1}
\]
The last equality identifies \(R\) with scalar multiplication on
the inclusion \(I_M\subset S\). Consequently every graded map
\(I_M\to\bigoplus_{s\geq N}S e_s\), with \(N\geq1\), is zero.

\begin{proof}
Rationally \(S\otimes_R\mathbb Q=\mathbb Q[T]\), and
\(I_M\otimes_R\mathbb Q=T^M\mathbb Q[T]\).
A weight-preserving map to \(\mathbb Q[T]e_s\) therefore has,
on every sufficiently high weight, the form
\[
 T^{[w]}\longmapsto
       \lambda\,\frac{(w-s)!}{w!}\,T^{[w-s]}e_s
             \quad(\lambda\in\mathbb Q).               \tag{136.NT2}
\]
This is multiplication by \(\lambda T^{-s}\) under the ordinary
polynomial identifications. If an integral map is nonzero, its
rationalization is nonzero because the target is \(R\)-torsion-free.
For \(s>0\), take \(w=p^a\geq\max(M,s)\). The denominator
\(w(w-1)\cdots(w-s+1)\) in (136.NT2) has \(p\)-valuation at
least \(a\). Its valuation is unbounded, whereas that of a fixed
nonzero \(\lambda\) is finite. Formula (136.NT2) cannot then have
all its coefficients in \(R\). This proves the first assertion.
For \(s=0\), the coefficient is just \(\lambda\), so it must lie
in \(R\), and every such scalar works.
Projection to each summand proves the direct-sum assertion; no
finite-generation assertion about \(I_M\) is used.
\end{proof}

\paragraph{ARCH-136-NATIVE-TAIL-NOSPLIT.}
Fix a prime \(p\), \(m\geq2\), and \(N\geq1\).
There exists an integral admissible theory \(C\), with
\(C_0=S_{\mathbb Z}\) and horizontal cutoffs \(\mu_C(q)=q+1\),
and a guarded pure horizontal source map \(H\to C\), with a
selected odd disk generator \(a\in H_1^{m,1}\), such that, writing
\[
       b=da,\qquad z=\gamma_p(b),\qquad
       e=a\gamma_{p-1}(b),\qquad de=pz,                 \tag{136.NT3}
\]
the following properties hold in \(C_1\):
\begin{enumerate}
\item the image of \([z]\) is nonzero, of order \(p\), in native
weight \(pm\);
\item for every \(s\geq N\), there are individual Moore-pair
primitives \(u_s,v_s\) with
\(du_s=T^{[s]}z\) and \(dv_s=pu_s-T^{[s]}e\);
\item for no \(M\geq1\) is there a weight-preserving
\(S_{\mathbb Z}\)-linear family of primitives on the whole tail:
\[
 h:I_M z\longrightarrow C_1^{*,2p-1},\qquad dh(u z)=u z
                          \quad(u\in I_M).
                                                        \tag{136.NT4}
\]
\end{enumerate}
Here \(I_Mz\) in (136.NT4) denotes the coefficient module freely
identified with \(I_M\), shifted by the weight \(pm\); its map
to \(C_1\) is the displayed coefficient action. In fact the
extra vertical cohomology of \(C\) is zero in every form degree
as soon as the weight is at least \(pm+N\).

\paragraph{The coefficient complex.}
We first work over \(R=\mathbb Z_{(p)}\), and put
\(S=\Gamma_R(T)\). Take the free graded cover
\[
 F=\bigoplus_{s\geq N}S e_s\xrightarrow{\epsilon}I_N,
 \qquad\operatorname{wt}(e_s)=s,\quad
 \epsilon(e_s)=T^{[s]},\qquad J=\ker\epsilon.
                                                        \tag{136.NT5}
\]
Let \(L\) be the cochain complex
\[
       L^{-2}=J\xrightarrow{\mathrm{incl}}L^{-1}=F
          \xrightarrow{\epsilon}L^0=S,                 \tag{136.NT6}
\]
where the last map includes \(I_N\) in \(S\).
It has only \(H^0(L)=S/I_N\), and the inclusion
\(S[0]\to L\) is a chain map. The negative degrees in this
coefficient complex are deliberate. The construction below tensors
them only with a top-arity complex of minimum form degree \(p\),
so that the resulting PD theory has no negative form degrees.
All weight components of \(L\) are finite free over \(R\):
this is clear for \(S,F\), and \(J^w\) is a submodule of a finite
free module over the principal ideal domain \(R\).
The module \(J\) is nevertheless not flat over \(S\).
Otherwise ARCH-136-INTEGRAL-FLAT-TAIL applied to the bounded
complex \(L\), whose only cohomology is the finite-weight module
\(S/I_N\), would make multiplication by \(p\) surjective on
\(S/I_N\), contrary to its weight-zero copy of \(R\).

\paragraph{A natural nonnegative PD construction.}
For a free \(R\)-module \(V\), put
\[
 D(V)=S\otimes_R\Lambda_R(aV)\otimes_R\Gamma_R(bV),
 \quad da_v=b_v,
 \quad |a_v|_{\rm form}=1,\quad |b_v|_{\rm form}=2,
                                                        \tag{136.NT7}
\]
giving both copies weight \(m\) and auxiliary arity one.
Let \(D_r(V)\) denote its arity-\(r\) part, including the
coefficient factor \(S\). Truncate all arities above \(p\).
In that truncated algebra its arity-\(p\) ideal is square-zero
and is annihilated by every positive arity.

Replace that top ideal by the complex
\[
       M(V)=D_p(V)\otimes_S L,
 \qquad D_p(V)\longrightarrow M(V)
                    \text{ induced by }S[0]\longrightarrow L.
                                                        \tag{136.NT8}
\]
Explicitly define the underlying complex of \(E(V)\) to be
\[
       S\ \oplus\!\bigoplus_{1\leq r<p}D_r(V)
                                  \ \oplus M(V).       \tag{136.NT9}
\]
Products of old terms of total arity less than \(p\) are their
old products. Products of total arity \(p\) are their old products
followed by the map in (136.NT8). Products of larger arity vanish.
The new top module has its usual \(S\)-action and is annihilated
by every positive arity. The same rules define the divided powers:
old operations are followed by the map to the top module when
their arity is \(p\), and operations of larger arity vanish;
all higher divided powers of a new top-ideal element are zero.
Mixed sums are defined by the divided-power addition formula.

These rules give a strict differential PD algebra.
Indeed every product or PD identity involving only old elements
is the corresponding identity in the truncated algebra, followed
by the map of its top ideal. If a new top element is involved,
every product with positive arity, and every product of two top
elements, is zero; the remaining identities are those of the
special square-zero PD module over \(S\). For example, on even
zero-forms the only extra formula is
\(\gamma_r(s+u)=\gamma_r(s)+\gamma_{r-1}(s)u\) for a
coefficient \(s\) and top element \(u\). The differential of
the top module is \(S\)-linear, and its inclusion map is a chain
map. Leibniz and the special identity
\(d\gamma_r(x)=\gamma_{r-1}(x)dx\) therefore follow from
the old identities and the same annihilation rules.

Since \(D_p(V)\) has form degrees between \(p\) and \(2p\),
\(M(V)\) has form degrees at least \(p-2\geq0\).
Thus \(E(V)\) is nonnegative, including at \(p=2\).
Every construction is natural for linear maps of \(V\), and
there is a natural strict augmentation \(E(V)\to S\).

\paragraph{Vertical cohomology and horizontal contraction.}
Let \(U\) be the simplicial free \(R\)-module corresponding
under Dold--Kan to the homological unit disk
\(R[2]\xrightarrow{1}R[1]\). It is simplicially contractible,
\(U_0=0\), and \(U_1=R\). Apply (136.NT9) simplexwise, and set
\[
                 C_R=B_R\otimes_S E(U).                \tag{136.NT10}
\]
This is an actual simplicial strict differential PD algebra.

For \(1\leq r<p\), the integral arity Euler operator
\(h(b_v)=a_v\), \(h(a_v)=0\), satisfies
\(dh+hd=r\) on \(D_r(V)\). Since \(r\) is a unit in \(R\),
these positive-arity complexes contract onto zero.
The terms of \(D_p(V)\) are free over \(S\), and this complex
is bounded. Tensoring the quasi-isomorphism
\(L\to(S/I_N)[0]\) therefore gives
\[
       M(V)\simeq D_p(V)\otimes_S(S/I_N).               \tag{136.NT11}
\]
The complex on the right is entirely supported in the weights
\(pm,pm+1,\ldots,pm+N-1\). Its forms are at least \(p\geq2\).
This proves that the coefficient map gives \(H^0(E(V))=S\),
and that every extra vertical cohomology group vanishes in
weight at least \(pm+N\), independently of the rank of \(V\).
The standard \(S\)-linear Poincare contraction of \(B_n\)
onto \(S\) remains a contraction after tensoring with any
complex. It identifies the vertical cohomology of \((C_R)_n\)
with that of \(E(U_n)\), with the same bound independent of \(n\).

Every positive-arity form row of \(E(U)\) is a finite direct
sum, in each fixed weight and form, of functors
\(D_r(U)\otimes_S W\), for constant coefficient modules
\(W\) among \(S,F,J\). Apply these functors to a linear
simplicial contraction of \(U\). Functoriality preserves all
simplicial homotopy identities, and the zero map acts as zero
in positive arity. Thus every such row is contractible.
The same holds after diagonal tensor with \(B_R\); explicitly
the homotopy uses \(s_i\) on the \(B_R\) factor and the
chosen \(i\)-th homotopy map on the other factor.
The augmentation splitting consequently identifies the horizontal
homology of \(C_R\) with that of \(B_R\). Its cutoff is
\(w>q\). Also \((C_R)_0=S\), because \(U_0=0\).

\paragraph{The native cycle and failure of every tail lift.}
At simplex one, \(U_1\) has rank one. The top arity block before
tensoring with \(L\) is exactly
\[
             S e\xrightarrow{p}S z
                     \quad\text{in forms }2p-1,2p,
                                                        \tag{136.NT12}
\]
with the notation (136.NT3). Formula (136.NT11) shows that its
new top cohomology is
\[
                H^{2p}(E(U_1))=(S/(pS+I_N))\,z.        \tag{136.NT13}
\]
In particular the native class is nonzero of order \(p\).
Each \(T^{[s]}z\), \(s\geq N\), is the differential of
\(z\otimes e_s\in D_p(U_1)^{2p}\otimes L^{-1}\).
More precisely put \(u_s=z\otimes e_s\) and
\(v_s=e\otimes e_s\). The total-complex sign on the odd
element \(e\) gives
\[
       du_s=T^{[s]}z,\qquad
       dv_s=pu_s-T^{[s]}e.
                                                        \tag{136.NT16}
\]
Thus the example retains the full individual Moore-pair lift,
not merely the vanishing of the first class.

Suppose there were an \(S\)-linear family (136.NT4) in
\((C_R)_1\). Evaluate its \(B_R\) factor at either vertex.
This is an \(S\)-linear chain map
\((C_R)_1\to E(U_1)\) and is the identity on the selected disk.
The form-\((2p-1)\) group of \(E(U_1)\) is precisely
\[
                         S e\ \oplus F z.              \tag{136.NT14}
\]
Lower arities have maximum form degree \(2p-2\), and the only
two summands of the tensor top block in that form are those shown.
After removing the common weight shift \(pm\), (136.NT1)
says that any map \(I_M\to S e\oplus Fz\) is
\(u\mapsto\lambda u e\), for one \(\lambda\in R\).
Its differential is \(p\lambda u z\). It cannot equal
\(u z\) for all \(u\in I_M\), since \(p\lambda=1\) has
no solution in \(R\). The coefficient action on the underlying
\(S z\) term is injective, so this comparison of coefficients
is legitimate. This proves failure of every tail lift over \(R\).

\paragraph{An integral theory and actual guarded source data.}
The preceding obstruction is not merely a statement about a
localized coefficient category. The strict augmentation
\(C_R\to B_R\) has its coefficient section. Form the strict
pullback of simplicial differential PD algebras
\[
           C=B_{\mathbb Z}\times_{B_R} C_R,
            \qquad B_{\mathbb Z}\longrightarrow B_R
                       \text{ the scalar localization}. 
                                                        \tag{136.NT15}
\]
As a complex of abelian groups this is
\(B_{\mathbb Z}\oplus\ker(C_R\to B_R)\), with the actual
coefficient inclusion from \(S_{\mathbb Z}\) in its first summand.
The second summand has the uniformly bounded-weight vertical
cohomology and the contractible horizontal rows just proved.
Hence \(C\) is an integral admissible theory, \(C_0=S_{\mathbb Z}\),
and \(\mu_C(q)=q+1\). It is a strict pullback, so all
divided-power and differential identities are retained.

The elements \(a,b,z,e\) above lie in the augmentation ideal,
so they define elements \((0,a),(0,b),(0,z),(0,e)\) in this
integral theory, with exactly the identities (136.NT3).
Its underlying extra summand is an \(R\)-module. Localizing
(136.NT15) at \(p\) therefore recovers \(C_R\): this follows
either from the displayed direct sum or from flatness of scalar
localization in the pullback exact sequence. Any hypothetical
integral family (136.NT4) would localize to the forbidden family
over \(R\). The native class and its individual tail boundaries
also survive in the integral direct summand. This proves all
three assertions for \(C\).

Finally the chosen element \(a\in C_1^{m,1}\) has zero
simplicial boundary, because the augmentation ideal is zero
in simplex zero. Since \(m\geq2=\mu_C(1)\), it can be the
image of the zero-boundary guarded horizontal cell in form one
and weight \(m\); its vertical partner maps to \(b\).
Continue the pure horizontal construction of
ARCH-136-DISK-SUPPORT. Acyclicity of every guarded target row
supplies every further chosen horizontal primitive, without
altering the initial cell map. This gives an actual guarded
source \(H\to C\), as asserted.

\paragraph{Exact consequence for nonflat tail repairs.}
A proposed adjunction that installs one copy of the coefficient
module \(I_Mz\) in form \(2p-1\), with its differential equal
to the coefficient multiples of \(z\), necessarily asks for
the map (136.NT4) in every target. This map need not exist even
for an integral admissible target, even though every individual
high-weight multiple is a boundary, and even though the starting
data occur in a genuine guarded free horizontal source.
Replacing that module by a free resolution returns to the
bounded-flat arity obstruction if the resolution stays bounded
in the natural arity.

The constructed target itself shows the distinction: it realizes
individual primitives through the free cover \(F\), and its
relations through the nonflat module \(J\), without choosing a
coefficient-linear primitive on \(I_M\). The proof does not rule
out other coherent nonflat presentations, source replacements,
or mixed-arity constructions. It does not decide the original
integral localization groupoid.
Indeed the constructed theory has a strict split augmentation
\(C\rightleftarrows B_{\mathbb Z}\), so it is already in the
standard object's connected component. The failure concerns the
prescribed guarded map and the specified coefficient-linear repair.
